\documentclass[11pt]{article}

\usepackage[margin=1in]{geometry}
\usepackage{amsmath}
\usepackage{amssymb}
\usepackage{hyperref}

\title{The Current \texttt{Dynamics} Class}
\author{}
\date{}

\begin{document}

\maketitle

\section{Purpose}

The current \texttt{Dynamics} class is the collision-response kernel for the
clean rebuild. It is intentionally separated from \texttt{Base}. The
\texttt{Base} class owns the simulation loop, particle storage, contact
detection, double-buffered position movement, and reporting. The
\texttt{Dynamics} class owns only the math that converts detected contacts into
momentum changes.

This separation is important because the dynamics code is the part most likely
to be moved later into C++ or GLSL. The current design is therefore written with
a simple kernel-like rule:

\[
    \text{one worker owns one particle}.
\]

During collision processing, the worker for particle \(i\) may read selected
state from another particle \(j\), but it writes only to particle \(i\).

\section{Particle Data Contract}

Each particle is currently stored as a Python dictionary. The dynamics class
expects the following fields to exist before collision processing begins.

\subsection{Position}

Position is double-buffered:

\[
\texttt{location}
=
\{
\texttt{use},\texttt{x1},\texttt{y1},\texttt{x2},\texttt{y2}
\}.
\]

If
\[
    \texttt{use}=1,
\]
then the active position is
\[
    \mathbf{x}=(x_1,y_1).
\]
If
\[
    \texttt{use}=2,
\]
then the active position is
\[
    \mathbf{x}=(x_2,y_2).
\]

The helper \texttt{current\_location()} reads this active position.

\subsection{Particle Properties}

The dynamics code reads:

\[
    r_i = \texttt{radius},
    \qquad
    m_i = \texttt{mass},
    \qquad
    \mathbf{v}_i = (v_{x,i},v_{y,i}).
\]

The current overlap-area formula assumes equal-radius disks:

\[
    r_i = r_j = r.
\]

If unequal radii are passed to \texttt{circle\_overlap\_area()}, the code raises
an error.

\subsection{Contact List}

\texttt{Base.detect\_contacts()} fills each particle's contact storage before
\texttt{Dynamics.process\_collisions()} runs. For particle \(i\), the relevant
fields are:

\[
    \texttt{contact\_count}_i,
    \qquad
    \texttt{contacts}_i.
\]

\texttt{contacts} is a fixed-length array of particle indices. Only the first
\texttt{contact\_count} entries are valid.

\section{Memory Ownership Rule}

The dynamics pass is particle-owned. When processing particle \(i\), the code
loops through \(i\)'s own contacts:

\[
    j \in \texttt{contacts}_i.
\]

The code may read particle \(j\)'s current location and radius to calculate
geometry. It does not write to particle \(j\). This is different from a
traditional pair solver that updates both particles in one pair operation.

The current update has the form:

\[
    \Delta \mathbf{p}_i
    =
    \sum_{j \in C_i}
    \Delta \mathbf{p}_{ij}.
\]

The corresponding update for particle \(j\) is produced only when particle
\(j\)'s own worker processes its own contact list.

\section{Contact Geometry}

For a pair \(i,j\), the active center positions are:

\[
    \mathbf{x}_i = (x_i,y_i),
    \qquad
    \mathbf{x}_j = (x_j,y_j).
\]

The separation vector from particle \(i\) to particle \(j\) is:

\[
    \mathbf{d}_{ij}
    =
    \mathbf{x}_j-\mathbf{x}_i.
\]

The center distance is:

\[
    d = \|\mathbf{d}_{ij}\|.
\]

Contact exists when:

\[
    d < r_i+r_j.
\]

If there is no overlap, \texttt{particle\_contact()} returns
\texttt{None}.

\subsection{Contact Normal}

When \(d>0\), the unit normal from particle \(i\) toward particle \(j\) is:

\[
    \mathbf{n}_{ij}
    =
    \frac{\mathbf{x}_j-\mathbf{x}_i}{d}
    =
    (n_x,n_y).
\]

If the centers are effectively identical,
\[
    d \le 10^{-12},
\]
the code uses the deterministic fallback:

\[
    \mathbf{n}_{ij}=(1,0).
\]

This avoids division by zero.

\section{Penetration Limit Used For Area}

The raw overlap depth is:

\[
    \delta_{\text{raw}}
    =
    r_i+r_j-d.
\]

The current code caps the depth used for area calculation:

\[
    \delta
    =
    \min\left(
        r_i+r_j-d,
        \min(r_i,r_j)
    \right).
\]

Under the equal-radius assumption, this means:

\[
    \delta \le r.
\]

This matches the current development assumption that particles do not move past
the case where the leading edge reaches the other particle center.

The distance sent to the overlap-area formula is:

\[
    d_A
    =
    r_i+r_j-\delta.
\]

\section{Equal-Radius Overlap Area}

For equal-radius particles,
\[
    r_i=r_j=r,
\]
the circular lens area is:

\[
    A(d_A)
    =
    2r^2\cos^{-1}\left(\frac{d_A}{2r}\right)
    -
    \frac{d_A}{2}\sqrt{4r^2-d_A^2}.
\]

The implementation clamps the distance into the valid range:

\[
    d_c
    =
    \min(\max(d_A,0),2r).
\]

The actual computed area is:

\[
    A(d_c)
    =
    2r^2\cos^{-1}\left(\frac{d_c}{2r}\right)
    -
    \frac{d_c}{2}\sqrt{\max(0,4r^2-d_c^2)}.
\]

The \(\max(0,\cdot)\) term protects the square root from small floating-point
roundoff errors.

\section{Overlap Momentum}

The current model follows the minimal legacy Mom-style idea:

\[
    \text{overlap area}
    \rightarrow
    \text{momentum amount}.
\]

It does not treat overlap as a force integrated over the substep. There is no
multiplication by \(\Delta t\) inside \texttt{overlap\_momentum()}.

For a contact between particle \(i\) and particle \(j\), the scalar momentum
amount is:

\[
    J_{ij}
    =
    k_i A_{ij} w(d).
\]

Here:

\[
    k_i = \texttt{momentum\_per\_area}
\]

and:

\[
    A_{ij} = \text{overlap area}.
\]

The coefficient \(k_i\) is owned by particle \(i\). If particle \(i\) has a
\texttt{momentum\_per\_area} field, that value is used. Otherwise the class
default is used. The code also keeps the legacy fallback name
\texttt{repulsion\_force\_per\_area}.

\section{Inverse-Square Weight}

The distance weight is:

\[
    w(d)
    =
    \frac{1}{\max(d^2,s^2)}.
\]

where

\[
    s = \texttt{inverse\_square\_softening}.
\]

The softening value prevents singular behavior when two centers are very close:

\[
    s \ge 10^{-12}.
\]

\section{Momentum Direction}

The contact normal \(\mathbf{n}_{ij}\) points from particle \(i\) toward
particle \(j\). The collision response for particle \(i\) points away from
particle \(j\), so the update for this contact is:

\[
    \Delta \mathbf{p}_{ij}
    =
    -J_{ij}\mathbf{n}_{ij}.
\]

Component-wise:

\[
    \Delta p_{x,ij} = -n_x J_{ij},
    \qquad
    \Delta p_{y,ij} = -n_y J_{ij}.
\]

\section{Accumulating Contacts}

For a particle with contact set \(C_i\), the dynamics class accumulates:

\[
    P_{x,i}
    =
    \sum_{j\in C_i}
    -n_{x,ij}J_{ij},
\]

\[
    P_{y,i}
    =
    \sum_{j\in C_i}
    -n_{y,ij}J_{ij}.
\]

It also tracks the scalar sum:

\[
    P_{\text{sum},i}
    =
    \sum_{j\in C_i}
    J_{ij}.
\]

These values are written back onto particle \(i\) as:

\[
\begin{aligned}
    \texttt{overlap\_momentum\_x} &= P_{x,i},\\
    \texttt{overlap\_momentum\_y} &= P_{y,i},\\
    \texttt{overlap\_momentum} &= P_{\text{sum},i}.
\end{aligned}
\]

The same vector is currently also stored as internal momentum:

\[
\begin{aligned}
    \texttt{internal\_momentum\_x} &= P_{x,i},\\
    \texttt{internal\_momentum\_y} &= P_{y,i},\\
    \texttt{internal\_momentum} &=
    \sqrt{P_{x,i}^2+P_{y,i}^2}.
\end{aligned}
\]

\section{Velocity Update}

The method \texttt{apply\_overlap\_momentum()} converts the stored momentum
vector into a velocity change:

\[
    \Delta \mathbf{v}_i
    =
    \frac{\Delta \mathbf{p}_i}{m_i}.
\]

Component-wise:

\[
    v_{x,i}^{\text{new}}
    =
    v_{x,i}
    +
    \frac{\texttt{overlap\_momentum\_x}}{m_i},
\]

\[
    v_{y,i}^{\text{new}}
    =
    v_{y,i}
    +
    \frac{\texttt{overlap\_momentum\_y}}{m_i}.
\]

The position update for the current substep happens in \texttt{Base} before
collision processing. Therefore this velocity change affects the next movement
step.

\section{Current Step Order}

At the current rebuild stage, the relevant order is:

\begin{enumerate}
    \item \texttt{Base} predicts the next position using current velocity.
    \item \texttt{Base.move()} flips the active location buffer.
    \item \texttt{Base.detect\_contacts()} fills each particle's contact list.
    \item \texttt{Dynamics.process\_collisions()} computes overlap momentum.
    \item \texttt{Dynamics.apply\_overlap\_momentum()} updates velocity.
    \item \texttt{Base.record\_step\_report()} stores per-particle values.
\end{enumerate}

\section{What This Class Does Not Do}

The current \texttt{Dynamics} class does not include bonding, field dynamics,
long-term internal momentum release, classical restitution coefficient \(e\),
or pair-store bookkeeping. It is only the minimal overlap-momentum collision
kernel.

It also does not update both particles inside a single pair operation. That is
intentional for the current particle-owned memory model.

\section{Future Porting Notes}

The most direct C++ or GLSL version would give each particle one worker. That
worker would:

\begin{enumerate}
    \item read particle \(i\)'s contact count and contact array,
    \item loop through the valid contact slots,
    \item read particle \(j\)'s current position and radius,
    \item compute \(A_{ij}\), \(J_{ij}\), and \(-J_{ij}\mathbf{n}_{ij}\),
    \item write only particle \(i\)'s momentum outputs and updated velocity.
\end{enumerate}

The current Python class is shaped to keep that translation straightforward.

\end{document}
